\vspace{0.5em}\noindent\textbf{Architecture}\par

\vspace{0.5em}\noindent\textbf{Objective}\par

This chapter explains the final AWS architecture and the main runtime
flows of the Internship Application Tracker.

\vspace{0.5em}\noindent\textbf{Scope}\par

This architecture is based on the final production evidence, the project
documentation notes, and the current application repository. I treat
runtime evidence as the source of truth when it differs from older
manifests. The most important current correction is that the frontend is
hosted by S3 and CloudFront, not by EKS.

\vspace{0.5em}\noindent\textbf{Architecture context}\par

{
\begingroup
\small
\setlength{\tabcolsep}{3pt}
\renewcommand{\arraystretch}{1.15}
\begin{longtable}{@{}
  >{\raggedright\arraybackslash}p{(\linewidth - 2\tabcolsep) * \real{0.5000}}
  >{\raggedright\arraybackslash}p{(\linewidth - 2\tabcolsep) * \real{0.5000}}@{}}
\toprule\noalign{}
\begin{minipage}[b]{\linewidth}\raggedright
Layer
\end{minipage} & \begin{minipage}[b]{\linewidth}\raggedright
Main components
\end{minipage} \\
\midrule\noalign{}
\endhead
\bottomrule\noalign{}
\endlastfoot
Edge & CloudFront distribution \texttt{EQIGYNECXDYL8} \\
Public routing & Application Load Balancer
\texttt{k8s-\allowbreak{}internshippublic-\allowbreak{}48101b50ad-\allowbreak{}85486086.\allowbreak{}ap-\allowbreak{}southeast-\allowbreak{}1.\allowbreak{}elb.\allowbreak{}amazonaws.\allowbreak{}com} \\
Kubernetes & EKS cluster \texttt{internship-\allowbreak{}prod}, namespace
\texttt{internship} \\
Workloads & \texttt{backend}, \texttt{chat-\allowbreak{}service},
\texttt{backend-\allowbreak{}outbox-\allowbreak{}dispatcher}, \texttt{backend-\allowbreak{}processing-\allowbreak{}worker},
\texttt{ai-\allowbreak{}service} \\
Data & RDS PostgreSQL, DynamoDB, ElastiCache Redis, S3 \\
Async and event & PostgreSQL processing jobs, PostgreSQL outbox, SQS,
Lambda, SES \\
AI & \texttt{ai-\allowbreak{}service} adapter and SageMaker endpoint
\texttt{internship-\allowbreak{}qwen3-\allowbreak{}4b} \\
CI/CD & GitHub Actions, AWS OIDC, ECR, kubectl, AWS CLI \\
\end{longtable}
\endgroup
}

\vspace{0.5em}\noindent\textbf{Diagram 1: Overall AWS
architecture}\par

\textbf{Flow summary.} Users enter through CloudFront, which serves the
private S3 frontend or forwards dynamic requests to the ALB and EKS. EKS
workloads connect to RDS, DynamoDB, Redis, S3, SageMaker, and the
SQS/Lambda notification path, with a DLQ and DynamoDB deduplication
table for resilient event handling.

\vspace{0.5em}\noindent\textbf{Diagram 2: CloudFront request-routing
flow}\par

\textbf{Flow summary.} CloudFront evaluates each HTTPS path. The default
behavior serves the S3 frontend; \texttt{/\allowbreak{}api/\allowbreak{}*} is forwarded to the
backend service; \texttt{/\allowbreak{}chat/\allowbreak{}*} and \texttt{/\allowbreak{}socket.\allowbreak{}io/\allowbreak{}*} are
forwarded to the chat service through the ALB.

The frontend bucket remains private. CloudFront handles HTTPS and SPA
fallback. The ALB Ingress rewrites \texttt{/\allowbreak{}api} and \texttt{/\allowbreak{}chat}
prefixes before forwarding to the backend and chat services.

\vspace{0.5em}\noindent\textbf{Diagram 3: CV upload and AI processing
sequence}\par

\textbf{Flow summary.} A candidate submits a CV through the frontend,
and the backend stores the document reference plus application,
idempotency, outbox, and processing records. The worker claims the
queued job, calls the AI adapter and SageMaker, saves the normalized
result in PostgreSQL, and exposes the result for frontend polling.

\vspace{0.5em}\noindent\textbf{Diagram 4: Realtime chat
sequence}\par

\textbf{Flow summary.} Socket.IO traffic passes from CloudFront and the
ALB to a chat-service pod. The receiving pod persists the message in
DynamoDB, publishes the room event through Redis so other pods can
notify subscribed users, and acknowledges the sender only after
persistence.

DynamoDB is the permanent chat store. Redis is only the cross-pod
pub/sub layer and does not replace persistent storage.

\vspace{0.5em}\noindent\textbf{Diagram 5: Transactional outbox, SQS, and Lambda
sequence}\par

\textbf{Flow summary.} The backend writes the business mutation and
outbox event in one PostgreSQL transaction. The dispatcher publishes the
event to SQS, after which Lambda uses a conditional DynamoDB write for
deduplication; first-time events are archived in S3 and may trigger SES,
while duplicates are treated as already processed.

SQS delivery is at least once. Duplicate delivery is expected, so Lambda
idempotency is required. The successful smoke-test event was
\texttt{lambda-\allowbreak{}smoke-\allowbreak{}fixed-\allowbreak{}1785220478}, archived under
\texttt{outbox-\allowbreak{}archive/\allowbreak{}2026/\allowbreak{}07/\allowbreak{}28/\allowbreak{}lambda-\allowbreak{}smoke-\allowbreak{}fixed-\allowbreak{}1785220478.\allowbreak{}json}.

\vspace{0.5em}\noindent\textbf{Diagram 6: Source-code change and CI/CD deployment
flow}\par

\textbf{Flow summary.} A manually triggered GitHub Actions workflow
validates the repository and restores compute when required, then builds
SHA-tagged ECR images. It deploys the EKS application and public ALB,
ensures the private S3/CloudFront configuration, publishes the frontend,
invalidates CloudFront, and produces a deployment summary.

The workflow supports \texttt{validate}, \texttt{deploy},
\texttt{rollout}, \texttt{restore-\allowbreak{}compute}, \texttt{deploy-\allowbreak{}app},
\texttt{deploy-\allowbreak{}public}, \texttt{deploy-\allowbreak{}frontend}, and \texttt{full}.
GitHub Actions assumes an AWS role through OIDC; long-lived AWS access
keys are not part of the deployment design.

\vspace{0.5em}\noindent\textbf{Component descriptions}\par

{
\begingroup
\small
\setlength{\tabcolsep}{3pt}
\renewcommand{\arraystretch}{1.15}
\begin{longtable}{@{}
  >{\raggedright\arraybackslash}p{(\linewidth - 4\tabcolsep) * \real{0.3333}}
  >{\raggedright\arraybackslash}p{(\linewidth - 4\tabcolsep) * \real{0.3333}}
  >{\raggedright\arraybackslash}p{(\linewidth - 4\tabcolsep) * \real{0.3333}}@{}}
\toprule\noalign{}
\begin{minipage}[b]{\linewidth}\raggedright
Component
\end{minipage} & \begin{minipage}[b]{\linewidth}\raggedright
Responsibility
\end{minipage} & \begin{minipage}[b]{\linewidth}\raggedright
Verification command
\end{minipage} \\
\midrule\noalign{}
\endhead
\bottomrule\noalign{}
\endlastfoot
CloudFront & Public HTTPS entry point and routing &
\texttt{aws\ cloudfront\ get-\allowbreak{}distribution\ -\/\allowbreak{}-id\ EQIGYNECXDYL8} \\
S3 frontend bucket & Static frontend assets &
\texttt{aws\ s3\ ls\ s3:/\allowbreak{}/\allowbreak{}internship-\allowbreak{}prod-\allowbreak{}frontend-\textless{}AWS\_\allowbreak{}ACCOUNT\_\allowbreak{}ID\textgreater{}/\allowbreak{}} \\
ALB & Routes API/chat/socket traffic to EKS services &
\texttt{kubectl\ get\ ingress\ internship-\allowbreak{}public\ -n\ internship} \\
Backend & FastAPI business API and health endpoints &
\texttt{kubectl\ rollout\ status\ deployment/\allowbreak{}backend\ -n\ internship} \\
Chat service & Socket.IO and chat REST API &
\texttt{kubectl\ rollout\ status\ deployment/\allowbreak{}chat-\allowbreak{}service\ -n\ internship} \\
Processing worker & Async AI/document jobs &
\texttt{kubectl\ get\ deployment\ backend-\allowbreak{}processing-\allowbreak{}worker\ -n\ internship} \\
Outbox dispatcher & Publishes PostgreSQL events to SQS &
\texttt{kubectl\ logs\ deployment/\allowbreak{}backend-\allowbreak{}outbox-\allowbreak{}dispatcher\ -n\ internship} \\
AI service & SageMaker adapter with stable worker routes &
\texttt{kubectl\ port-\allowbreak{}forward\ service/\allowbreak{}ai-\allowbreak{}service\ 8010:8010\ -n\ internship} \\
RDS PostgreSQL & Transactional data and queues &
\texttt{aws\ rds\ describe-\allowbreak{}db-\allowbreak{}instances\ -\/\allowbreak{}-db-\allowbreak{}instance-\allowbreak{}identifier\ internship-\allowbreak{}prod-\allowbreak{}postgres\ -\/\allowbreak{}-region\ ap-\allowbreak{}southeast-\allowbreak{}1} \\
DynamoDB & Chat and Lambda dedupe tables &
\texttt{aws\ dynamodb\ describe-\allowbreak{}table\ -\/\allowbreak{}-table-\allowbreak{}name\ ChatMessages\ -\/\allowbreak{}-region\ ap-\allowbreak{}southeast-\allowbreak{}1} \\
Redis & Socket.IO pub/sub &
\texttt{aws\ elasticache\ describe-\allowbreak{}replication-\allowbreak{}groups\ -\/\allowbreak{}-replication-\allowbreak{}group-\allowbreak{}id\ internship-\allowbreak{}prod-\allowbreak{}redis\ -\/\allowbreak{}-region\ ap-\allowbreak{}southeast-\allowbreak{}1} \\
SQS & Outbox event transport &
\texttt{aws\ sqs\ get-\allowbreak{}queue-\allowbreak{}attributes\ -\/\allowbreak{}-queue-\allowbreak{}url\ \textless{}OUTBOX\_\allowbreak{}QUEUE\_\allowbreak{}URL\textgreater{}\ -\/\allowbreak{}-attribute-\allowbreak{}names\ All} \\
\end{longtable}
\endgroup
}

\vspace{0.5em}\noindent\textbf{Public and private
resources}\par

{
\begingroup
\small
\setlength{\tabcolsep}{3pt}
\renewcommand{\arraystretch}{1.15}
\begin{longtable}{@{}
  >{\raggedright\arraybackslash}p{(\linewidth - 2\tabcolsep) * \real{0.5000}}
  >{\raggedright\arraybackslash}p{(\linewidth - 2\tabcolsep) * \real{0.5000}}@{}}
\toprule\noalign{}
\begin{minipage}[b]{\linewidth}\raggedright
Public-facing
\end{minipage} & \begin{minipage}[b]{\linewidth}\raggedright
Private or internal
\end{minipage} \\
\midrule\noalign{}
\endhead
\bottomrule\noalign{}
\endlastfoot
CloudFront domain & RDS PostgreSQL \\
ALB DNS origin & ElastiCache Redis \\
CloudFront HTTPS routes & DynamoDB tables \\
Browser-accessible static assets through CloudFront & EKS pod-to-service
traffic \\
SES email delivery & SQS queue and DLQ \\
& S3 buckets without direct public bucket access \\
\end{longtable}
\endgroup
}

\vspace{0.5em}\noindent\textbf{Security boundaries}\par

\begin{itemize}
\tightlist
\item
  GitHub Actions uses OIDC to assume \texttt{internship-\allowbreak{}github-\allowbreak{}deploy}.
\item
  EKS workloads use \texttt{internship-\allowbreak{}app} ServiceAccount and IRSA role
  mapping.
\item
  Secrets are injected through GitHub environment secrets and Kubernetes
  \texttt{internship-\allowbreak{}secrets}.
\item
  S3 frontend bucket should remain private and readable through
  CloudFront only.
\item
  Backend and chat pods expose only service ports through the ALB path
  rules.
\item
  AI traffic stays inside the cluster from worker to
  \texttt{ai-\allowbreak{}service}; SageMaker is invoked by the adapter.
\item
  SQS/Lambda processing is idempotent through DynamoDB conditional
  writes.
\end{itemize}

\vspace{0.5em}\noindent\textbf{Expected result}\par

The final architecture separates static delivery, public request
routing, long-running service execution, transactional storage, realtime
synchronization, asynchronous event delivery, short event-driven
notification, and AI inference. Each component has a clear
responsibility and an independent verification path.

\vspace{0.5em}\noindent\textbf{Common errors}\par

{
\begingroup
\small
\setlength{\tabcolsep}{3pt}
\renewcommand{\arraystretch}{1.15}
\begin{longtable}{@{}
  >{\raggedright\arraybackslash}p{(\linewidth - 4\tabcolsep) * \real{0.3333}}
  >{\raggedright\arraybackslash}p{(\linewidth - 4\tabcolsep) * \real{0.3333}}
  >{\raggedright\arraybackslash}p{(\linewidth - 4\tabcolsep) * \real{0.3333}}@{}}
\toprule\noalign{}
\begin{minipage}[b]{\linewidth}\raggedright
Error
\end{minipage} & \begin{minipage}[b]{\linewidth}\raggedright
Cause
\end{minipage} & \begin{minipage}[b]{\linewidth}\raggedright
Resolution
\end{minipage} \\
\midrule\noalign{}
\endhead
\bottomrule\noalign{}
\endlastfoot
Frontend shown as running in EKS & Old architecture description & Use
current S3 and CloudFront architecture only \\
S3 described behind ALB & Incorrect CloudFront origin model & Default
CloudFront behavior must route directly to S3 \\
Redis described as chat database & Misunderstood chat design & DynamoDB
stores chat records; Redis is pub/sub \\
Outbox described as exactly-once & SQS Standard semantics & Document
at-least-once plus idempotent consumer \\
Worker enabled before AI readiness & SageMaker dependency not ready &
Keep worker disabled until AI and endpoint are ready \\
\end{longtable}
\endgroup
}

\vspace{0.5em}\noindent\textbf{Outcome}\par

This architecture chapter provides the reference model for the remaining
workshop steps. All later deployment, testing, monitoring, security,
cost, troubleshooting, and cleanup sections should match these diagrams.
